\chapter{Instance Fusion: Analysis and Improvements}
\label{ch:fusion}

% PRESUPUESTO: ~7 pp. UNIFICA el viejo diagnóstico de fusión + la revisión del criterio.
% Aquí NO nace el contest (esa premisa, el scope limit, vive ya en \ref{ch:contest}).
% Este capítulo diseca el merge POR DENTRO y lo mejora, sobre terreno firme (banco +
% grading + todas las mejoras agnósticas ya activas). Es la redención de la fase perdida:
% el mismo estudio de heurísticas del principio, ahora ATRIBUIBLE.

\section{Method: grading fusion in isolation}
\label{sec:fus_method}
% La maquinaria específica de fusión (movida aquí desde el framework): sirve solo para
% puntuar decisiones de fusión, así que aterriza donde se usa.

\subsection{The synthetic loop-closure protocol}
\label{sec:fus_protocol}
% Salto -> recorrido -> corrección oráculo con GT (trayectoria Y puntos) = LC sintético.
%   - corrección oráculo: el fallo es atribuible al mecanismo semántico, no al SLAM
%   - duplicados por construcción: se sabe qué par es duplicado de qué -> grading posible
%   - tres poblaciones: pre-drift | post-drift (contaminadas) | post-corrección (solapadas)
% LIMITACIÓN, escrita por nosotros: corrección perfecta, instantánea, única. Un LC real
% corrige imperfecto, progresivo, solo sobre el subgrafo de covisibilidad. El sintético
% AÍSLA el mecanismo; ORB-SLAM2 demuestra que sobrevive al ruido real. (Se apoya en
% SimulatedSLAM, \ref{sec:bench_simslam}.)

\subsection{Checkpoints: replaying fusion in isolation}
\label{sec:fus_checkpoints}
% Checkpoint pre-fusión -> replay desde estado idéntico cambiando SOLO la heurística.
% Convierte una comparación en un experimento controlado.

\subsection{Grading fusion decisions against ground truth}
\label{sec:fus_grading}
% fusion_metrics: matriz de confusión (TP/FP/FN/TN) de las decisiones frente al GT.
% Segmentaciones: global, pre-drift vs post-drift, etc.

\section{Where the merge spends its effort}
\label{sec:fus_scope}
% Diagnóstico del merge por dentro (viejo cap. failure-modes, sin el scope limit).

\subsection{Cost and selectivity are inverted}
\label{sec:fm_cost_selectivity}
% Centroide y coseno filtran muchísimo y son baratos. Overlap es el criterio MÁS caro y
% el que MENOS pares llega a valorar. El orden del pipeline está mal (se corrige en
% \ref{sec:assoc_ordering}). TABLA: coste x pares evaluados x pares rechazados por criterio.

\subsection{The true-negative flood}
\label{sec:fm_tn_flood}
% Enorme cantidad de TN -> las métricas globales no discriminan. Dos causas: comparación
% fuerza bruta todos-contra-todos, y sobre-segmentación. -> motiva la covisibilidad.

\section{Candidate selection by covisibility}
\label{sec:assoc_candidates}
% Covisibilidad en lugar de fuerza bruta todos-contra-todos. Doble pago: mata la
% inundación de TN (\ref{sec:fm_tn_flood}) -> métricas informativas, y recorta el coste
% -> la eficiencia que promete la propuesta oficial.
% RIESGO: puede no llegar a implementarse -> entonces se va a \ref{sec:concl_future} con
% su motivo, y el capítulo se sostiene con las otras secciones.

\section{Criterion ordering by cost and selectivity}
\label{sec:assoc_ordering}
% Reordenar el pipeline según lo medido en \ref{sec:fm_cost_selectivity}: barato y
% selectivo primero (centroide, coseno), caro al final (overlap). Ahora un cambio de
% umbral es ATRIBUIBLE a un cambio de resultado.

\section{Semantic descriptors}
\label{sec:assoc_descriptors}
% La única veta AWARE del capítulo. PE-Core (mejor), PE-Spatial, SigLip. Bounding box con
% y sin fondo. Qué aporta el descriptor y qué NO: no era el cuello de botella.
% Descarte con su motivo: distancia entre bounding boxes -> sin impacto.

\section{Results}
\label{sec:fus_results}
% Banco sintético (\ref{ch:fusion}, método) + ORB-SLAM2 real. Todas las mejoras activas.
